%% FeatureLiftBench: evidence-grounded manuscript revision, 2026-09-07.
%% Sources: python150_prime_v2_analysis_20260905/task_results.csv;
%% python150_paper_analysis_final/{summary_final.md,json/stats.json};
%% artifacts/research_analysis/python200_prime/current_benchmark_freeze.json.
%% Main scores retain the frozen 150-task denominator. No new agent runs.
\documentclass[acmsmall,screen,review,anonymous]{acmart}
\usepackage{booktabs}
\usepackage{tabularx}
\usepackage{graphicx}
\usepackage{xspace}
\usepackage{microtype}
\usepackage{placeins}
\newcommand{\flb}{\textsc{FeatureLiftBench}\xspace}
\newcommand{\rres}{\textsc{RRES}\xspace}
\newcommand{\mainprotocol}{\textsc{Main}\xspace}
\graphicspath{{figures/}}
\setcopyright{none}
\settopmatter{printacmref=false}
\acmDOI{}
\acmISBN{}
\acmPrice{}
\acmYear{2026}
\copyrightyear{2026}
\citestyle{acmauthoryear}
\begin{document}
\title{FeatureLiftBench: Evaluating Behavior-Preserving Feature Extraction from Real Repositories}
\author{Anonymous Author(s)}

\begin{abstract}
Reusing an existing software capability requires more than finding its main
implementation. Helpers, state, exception behavior, and resources must remain
consistent when the capability moves behind a new package boundary. We introduce
\flb, a benchmark that evaluates this process as \emph{feature lifting}:
an agent receives a pinned source repository and a public behavioral contract,
then constructs a standalone artifact that is evaluated without access to the
source project. The frozen release contains 200 Python tasks. Our main study
compares six model backends in OpenHands on a common 150-task subset, with source
location hints and benchmark tests withheld. Functional pass rates range from
24.0\% to 76.7\%. Failures concentrate at behavioral evaluation: only four of
829 delivered artifacts pass the Build, Public, and Hidden gates but fail the
separate Isolation gate. The in-suite hard3 construction group is associated
with lower success even after controlling for model backend and lift type.
Correctness and extraction footprint also differ: on 97 tasks passed by both
DeepSeek V4 Pro and GPT-5.6 Luna, median reference-relative size is 0.993 and
0.697, respectively. We separate these measured outcomes from exploratory
AI-assisted failure interpretation, disclose contract defects flagged during
analysis, and report a sensitivity view over existing results. \flb provides
an operational setting for studying both preservation of required behavior and
the construction of independently executable reusable artifacts.
\end{abstract}

\begin{CCSXML}
<ccs2012>
 <concept><concept_id>10011007.10011074.10011099.10011102</concept_id>
 <concept_desc>Software and its engineering~Software testing and debugging</concept_desc>
 <concept_significance>500</concept_significance></concept>
 <concept><concept_id>10011007.10011074.10011099</concept_id>
 <concept_desc>Software and its engineering~Software verification and validation</concept_desc>
 <concept_significance>300</concept_significance></concept>
</ccs2012>
\end{CCSXML}
\ccsdesc[500]{Software and its engineering~Software testing and debugging}
\ccsdesc[300]{Software and its engineering~Software verification and validation}
\keywords{coding agents, software reuse, feature extraction, behavioral preservation, benchmark}
\maketitle

\section{Introduction}
\label{sec:introduction}

A useful capability often exists inside a system whose surrounding dependencies
are unnecessary for a new application. A developer may want a parser without
its command-line frontend, a configuration resolver without its host framework,
or a plugin registry without the rest of a service. Moving that capability
requires deciding what must cross the new boundary. Copying an entry function
can omit transitive helpers, shared state, exception classes, or resource files;
copying a large part of the project can preserve behavior while retaining a
substantial implementation footprint. The engineering objective is therefore
twofold: preserve the requested behavior and construct a usable independent
artifact.

Coding agents can search repositories, inspect tests and documentation, edit
multiple files, and execute validation commands. These tools make feature
extraction a plausible application, but familiar benchmark scores do not directly
measure the complete operation. SWE-bench evaluates issue resolution within an
existing project~\cite{jimenez2024swebench}; HumanEval evaluates code generated
from function specifications~\cite{chen2021codex}. Feature location studies which
source entities implement a concern~\cite{dit2013featurelocation}. Software
transplantation already investigates extracting and adapting functionality for
reuse~\cite{barr2015transplantation}. Our contribution builds on this reuse
objective by defining a controlled, executable evaluation for repository-level
coding agents.

We call the task \emph{behavior-preserving feature lifting}. The agent is given
the complete source repository, including the original implementation, together
with a contract describing the required output behavior. It produces a new
\texttt{featurelifted} package, which is evaluated in a source-free environment.
The distinguishing boundary is the conjunction of source-backed behavior
recovery and independent execution of the delivered artifact. Neither a file
location nor a patch that relies on the original project is sufficient. A
behaviorally equivalent rewrite is permitted; copying a particular implementation
is not required.

\flb operationalizes this task with pinned Python repositories, explicit API and
behavior clauses, protected evaluation, and separate extraction-size metrics.
Under the Full-Repository / No-Hint \mainprotocol protocol, the agent can inspect
upstream code, tests, documentation, and resources, but cannot see benchmark
tests, source-entrypoint hints, evaluator internals, or reference solutions.
The evaluator records four gates: Build, Public, Hidden, and Isolation. The
Public and Hidden names distinguish evaluator test groups; both groups are
withheld during \mainprotocol. A functional pass requires all four gates.
After evaluation, trusted static analysis measures the submitted footprint
relative to a frozen feasible reference and estimates direct source copying.

The release contains 200 tasks from 176 canonical source repositories. The main
study uses the common Python-150 subset evaluated by six backends in OpenHands.
It includes an existing 100-task construction group and a 50-task hard3 group;
the additional 50 tasks in the release are a separate extension. This distinction
matters because these two 50-task groups have different roles. We report the
extension separately and do not use it to support the main difficulty analysis.

The study asks four questions. \textbf{RQ1:} How often do the evaluated
model--harness configurations produce functionally passing artifacts?
\textbf{RQ2:} Where do submissions fail in the evaluator?
\textbf{RQ3:} How does observed difficulty vary across tasks and construction
groups? \textbf{RQ4:} On common passing tasks, how do extraction size and direct
copying differ across backends?

The results show substantial but incomplete success. Pass rates span
24.0\%--76.7\%, and strong backends fail primarily on required behavior after
producing a loadable package. The hard3 group remains associated with lower
success in an analysis controlling for backend and lift type. Meanwhile,
functionally passing artifacts can have markedly different footprints on the
same tasks. An exploratory close-read of strong-backend failures suggests
incomplete behavioral-contract recovery as one recurring explanation. We keep
this interpretation distinct from the measured gate outcomes: the labels are
AI-assisted, and the analysis also flagged potential benchmark defects.

The paper makes three contributions:
\begin{itemize}
 \item \textbf{A feature-lifting task and benchmark.} We define an executable
 source-to-artifact task and assemble a frozen 200-task Python release with
 explicit output contracts and pinned source identities.
 \item \textbf{An evaluation protocol.} We specify the agent's information
 boundary, source-free functional evaluation, and pass-conditioned footprint
 measurement, with auditable task, image, and result records.
 \item \textbf{An empirical characterization.} We analyze six backends on
 150 common tasks, separating capability, evaluator failures, construction-group
 associations, paired extraction statistics, and exploratory failure evidence.
\end{itemize}

% RELATED_WORK_INSERT

\section{The FeatureLiftBench Benchmark}
\label{sec:benchmark}

\subsection{Task and Output Boundary}

A task consists of a pinned repository $R$, a public contract $S$, an initial
workspace $W$, and a protected evaluator $J$. A model backend $M$ acting through
an agent harness $H$ produces a submission $A$:
\begin{equation}
 A = \operatorname{Run}(M,H,R,S,W).
\end{equation}
The public contract lists the required API surface, behavior clauses, constraints,
and exclusions. Source identity records establish which upstream snapshot defines
the behavior being lifted. The agent can inspect that snapshot but must place the
result inside the designated submission directory.

The contract determines the scope of preservation. A task may require a selected
configuration operation or registry behavior without requiring the entire
upstream project. APIs, return values, exceptions, mutation, ordering, or resources
are obligations when specified by the task. Direct extraction, adapted code, and
reimplementation are legal strategies. This permits multiple passing artifacts
and avoids grading against a single required source slice.

The functional score is
\begin{equation}
 F(A)=B(A)\land T_{\mathrm{public}}(A)\land T_{\mathrm{hidden}}(A)\land I(A),
 \label{eq:functional}
\end{equation}
where $B$ is the evaluator's package loading/build gate, $T$ denotes the two
behavioral test groups, and $I$ denotes its explicit isolation checks. The tests
execute in an environment without the source repository. Consequently, a
dependency on a missing source path can fail during loading or behavioral tests,
before the separate Isolation gate becomes the first failing gate. Gate names
must not be interpreted as independent semantic failure causes.

\begin{figure}[t]
 \centering
 \fbox{\begin{minipage}{0.94\linewidth}
 \centering\small
 \textbf{Agent workspace}\\
 Public contract + complete pinned repository\\
 Upstream tests/docs available; benchmark tests and references withheld\\[5pt]
 $\downarrow$\quad OpenHands exploration, editing, and self-testing\quad$\downarrow$\\[5pt]
 \textbf{Submission boundary:}\; standalone \texttt{featurelifted} artifact\\[5pt]
 $\downarrow$\quad collect artifact only\quad$\downarrow$\\[5pt]
 \textbf{Source-free evaluation capsule}\\
 Build\quad$\land$\quad Public\quad$\land$\quad Hidden\quad$\land$\quad Isolation\\[5pt]
 $\downarrow$\\[3pt]
 Functional result + trusted, static extraction-size measurement
 \end{minipage}}
 \caption{Information and execution boundaries. The agent may inspect upstream
 tests, but both benchmark test groups remain evaluator-only in Main.}
 \Description{A source repository and public contract are available to the agent.
 Only the new artifact enters the source-free evaluation capsule. Four gates
 determine functional success, followed by static extraction measurements.}
 \label{fig:pipeline}
\end{figure}

\subsection{Task Construction and Freezing}
\label{sec:construction}

Construction links repository provenance, a public specification, protected
behavioral checks, and executable reference evidence. Candidate capabilities
come from Python libraries, developer tools, and framework components. The
selection targets capabilities that can be given a bounded contract and evaluated
offline. It is a curated collection, rather than a random sample of repositories
or a measure of the frequency of extraction requests in practice.

For each task, structured metadata records required APIs, behavioral clauses,
dependency information, and source provenance. The human-readable task statement
is rendered from the public specification. Evaluator-side records associate tests
with contract clauses, while the source registry records immutable revisions and
archive digests. These records allow changes in a task statement, evaluator, or
source snapshot to be distinguished from changes in an agent submission.

The freeze v2 repair process addressed undeclared API members, incorrect source
entrypoint metadata, and identical Public/Hidden test bodies. The repair scope
covered 38 distinct tasks. The records describe AI-assisted review and maintainer
adjudication of that repair scope; they do not constitute independent human
validation of every task. The frozen release records 200 successful task checks,
200 source mappings, and three passing oracle executions per task, for 600/600
executions with stable fingerprints (Table~\ref{tab:validation}). This establishes
executable feasibility and repeatability for the recorded references, rather
than completeness or fairness of every behavioral requirement.

\begin{table}[t]
 \caption{Recorded freeze v2 validation evidence and its scope.}
 \label{tab:validation}
 \centering\small
 \begin{tabularx}{\linewidth}{@{}lrlX@{}}
 \toprule
 Check & Coverage & Outcome & Interpretation \\
 \midrule
 Task validation & 200 tasks & Pass & Package/specification checks \\
 Source mapping & 200 tasks & Pass & Pinned source identity \\
 Oracle replay & 600 runs & Pass & Three executions per task \\
 Stable oracle results & 200 tasks & Stable & Recorded fingerprints agree \\
 Repair-scope review & 38 tasks & Preserved & AI-assisted, maintainer-adjudicated \\
 \bottomrule
 \end{tabularx}
\end{table}

We distinguish a \emph{validation rule} from a \emph{completed semantic audit}.
A protected expectation should be licensed by the public contract and supported
by upstream evidence. Mechanical API, entrypoint, hash, overlap, and execution
checks test parts of this relation; they cannot establish semantic entailment.
The release has not undergone a complete independent human audit of hidden-test
fairness. Later failure inspection flagged seven potential contract/evaluator
defects, which we disclose and use in a post hoc sensitivity analysis
(Section~\ref{sec:defects}). We do not silently repair the frozen tests or
reinterpret their recorded outcomes.

\subsection{Suite Composition}
\label{sec:suite}

The 200-task release consists of Python-150 and an additional 50-task extension
called Hard-50 in the artifact. The main six-backend comparison uses Python-150:
all six have results on these task IDs, and all passing submissions have the
reference-relative measurements used in this paper. Five backends also have
200-task results; Pro was evaluated on the 150-task subset. The 200-task
functional table is provided in Appendix~\ref{app:extension}, rather than
combining unequal model coverage into the main table.

Within Python-150, Core-100 and hard3-50 are construction groups. The
\emph{hard3-50 group belongs to Python-150}; it is not the additional Hard-50
extension. We retain these historical group names for reproducibility and test
their empirical relationship to performance without treating them as universal
categories of software difficulty.

Each task also has a lift-type label. \emph{Direct} denotes a target that can
largely retain its upstream structure; \emph{Adapted} denotes interface or
boundary changes; \emph{Composite} denotes behavior assembled from multiple
source regions or abstractions. Table~\ref{tab:construction} shows their
distribution. Feature-family and entanglement labels provide additional
descriptive metadata. They are hidden from the agent and do not affect scoring.
All headline tasks have the same high entanglement-level label, so that label
cannot support an empirical difficulty contrast here.

\begin{table}[t]
 \caption{Construction groups and lift types in the headline Python-150 subset.
 The additional Hard-50 release extension is outside this table.}
 \label{tab:construction}
 \centering
 \begin{tabular}{@{}lrrrr@{}}
 \toprule
 Group & Direct & Adapted & Composite & Total \\
 \midrule
 Core-100 & 53 & 44 & 3 & 100 \\
 hard3-50 & 3 & 32 & 15 & 50 \\
 \midrule
 Total & 56 & 76 & 18 & 150 \\
 \bottomrule
 \end{tabular}
\end{table}

\section{Evaluation Protocol and Study Design}
\label{sec:protocol}

\subsection{Agent Conditions and Recorded Runs}

We evaluate DeepSeek V4 Pro, DeepSeek V4 Flash, GPT-5.6 Luna served through
OpenLux, GLM-5.3-Flash, Qwen3.6-35B-A3B-FP8, and GPT-OSS 120B. Below, we
abbreviate them as Pro, Flash, Luna, GLM, Qwen, and OSS. The agent image records
OpenHands CLI 1.16.0. The common protocol supplies the complete pinned repository,
withholds benchmark tests and source-location hints, and uses the standard task
prompt. Agents can inspect upstream tests and create their own probes. Public
feedback and source-hint ablations are not pooled into this study.

The configured envelope is 120 steps, 131,072 context tokens with 8,192 reserved,
and a 3,600-second task timeout. Main does not impose the separate two-million
token cap used by the historical cost arm. Runtime statistics such as recorded
assistant messages are not necessarily identical to the configured action budget.
The evaluated object is a backend embedded in its recorded OpenHands/provider
profile, not a base model independent of tool calling, serving, and recovery.

The analysis retains one collected outcome per model--task cell. The campaign
includes replacements of preflight-blocked tasks and runner-level transient
recovery, recorded in the run artifacts. We therefore use Functional Pass@1 to
denote the final campaign outcome under this policy, not a claim that every task
made exactly one uninterrupted process invocation. There is no best-of-$k$
selection over independently scored candidate solutions in the reported analysis.
The final Python-150 matrix has 900 model--task cells. Empty submissions remain
failures in the 150-task denominator, including tool-validation and early-exit
events. Runner completion status is kept separate from evaluator correctness.

\subsection{Source-Free Functional Evaluation}

Only the collected submission enters the evaluator capsule. The original
repository is absent from the runtime import path, evaluation uses the frozen
dependency environment, and the container disables networking. The frozen image
records Linux/amd64 and Python 3.11.14, with a read-only root filesystem and
dropped capabilities. Agent and evaluator image identities are recorded in the
freeze manifest.

\textbf{Build} resolves and loads the submitted package according to the
evaluator's installation/path-loading policy. In this release, a direct package
directory can be loaded through \texttt{PYTHONPATH}; an editable-install failure
can also use a path fallback. Thus a Build pass demonstrates usability under the
frozen evaluator policy, not that a wheel or clean package-manager installation
has been independently validated. \textbf{Public} tests exercise primary
contract behavior, and \textbf{Hidden} tests cover additional cases and
combinations intended to follow the same public contract. \textbf{Isolation}
checks include forbidden source imports and protected-path restrictions.
All four outcomes are required by Equation~\ref{eq:functional}.

We report two complementary views. An exclusive first-outcome classification
uses Missing, Build, Public, Hidden, then Isolation, with Pass for a fully
successful artifact. A non-exclusive view counts failed gate flags among
delivered artifacts; these counts can overlap. A failed flag can reflect a
preceding loading failure and is not necessarily evidence that all tests in that
group executed. \emph{Isolation residual} denotes artifacts that pass Build,
Public, and Hidden but fail Isolation. This is a measured residual category,
not the prevalence of all source-dependency mistakes.

\subsection{Extraction Measurements}
\label{sec:metrics}

For a passing artifact $A$ and a frozen feasible reference $A_{\mathrm{ref}}$,
Reference-Relative Extraction Size is
\begin{equation}
 \operatorname{RRES}(A)=
 \frac{\operatorname{PythonLOC}(A)}{\operatorname{PythonLOC}(A_{\mathrm{ref}})}.
 \label{eq:rres}
\end{equation}
The evaluator uses its Python source-line counting convention consistently for
both artifacts. An RRES of one means equal counted size, not equal code or
optimality. The reference demonstrates one feasible implementation; it is not a
proof of a minimum semantic closure.

The copied-code diagnostic estimates the fraction of submitted lines matched
to source code in runs of at least three normalized non-empty, non-comment
Python lines. It can miss rewritten equivalents and count shared idioms. It is
neither a licensing judgment nor a direct measure of maintainability. These
metrics are computed in a trusted read-only stage after functional evaluation;
the metrics stage does not execute the submitted code.

Footprint summaries are conditioned on functional success. Because different
backends pass different tasks, unpaired medians are descriptive only. Comparisons
between backends use their same-task common passing set. Low RRES or low detected
copying cannot compensate for a functional failure, and neither is folded into
the primary pass score.

\subsection{Statistical and Evidence Boundaries}

We report Wilson 95\% intervals for functional pass proportions and selected
two-sided exact McNemar tests for paired model outcomes. Reported $p$-values are
unadjusted exploratory comparisons; they do not establish a total ordering of
backends. To study difficulty, we fit pass/fail against backend, construction
group, and lift type, with standard errors clustered by task. The model estimates
associations within this curated collection, not causal effects of construction.
For paired copying we report Wilcoxon signed-rank tests and matched-pairs
rank-biserial effect sizes.

One retained observation per cell does not estimate stochastic run-to-run
variance. Multiple tasks from a repository can also share behavior and failures;
task-level intervals and standard errors do not model every repository-level
dependence. The scores below retain the frozen task set. The separately labeled
known-defect sensitivity view is post hoc and uses only existing outcomes.

\section{Results}
\label{sec:results}

\subsection{RQ1: Functional Capability}
\label{sec:rq1}

Table~\ref{tab:main} reports results for all 150 assigned tasks per backend.
Pro has the highest observed pass rate at 115/150 (76.7\%), followed by Flash
at 108/150 (72.0\%) and Luna at 102/150 (68.0\%). GLM, Qwen, and OSS pass
68, 63, and 36 tasks, respectively. The observed range is 52.7 percentage
points. These scores include unsuccessful delivery, rather than conditioning
the denominator on whether the agent produced a package.

\begin{table}[t]
 \caption{Frozen Python-150 functional results. Empty submissions remain
 failures. Core and hard3 are the two groups inside this 150-task subset.}
 \label{tab:main}
 \centering\small
 \setlength{\tabcolsep}{4pt}
 \begin{tabular}{@{}lrrrrr@{}}
 \toprule
 Backend & Pass / 150 & Rate [Wilson 95\%] & Core / 100 & hard3 / 50 & Empty \\
 \midrule
 Pro & \textbf{115} & 76.7 [69.3, 82.7] & 94 & 21 & 0 \\
 Flash & 108 & 72.0 [64.3, 78.6] & 90 & 18 & 0 \\
 Luna & 102 & 68.0 [60.2, 74.9] & 82 & 20 & 6 \\
 GLM & 68 & 45.3 [37.6, 53.3] & 62 & 6 & 38 \\
 Qwen & 63 & 42.0 [34.4, 50.0] & 55 & 8 & 25 \\
 OSS & 36 & 24.0 [17.9, 31.4] & 25 & 11 & 2 \\
 \bottomrule
 \end{tabular}
\end{table}

The point estimates do not support every adjacent ranking. Pro succeeds alone
on ten tasks relative to Flash, while Flash succeeds alone on three; the exact
McNemar $p$-value is 0.0923. GLM versus Qwen has 27 versus 22 exclusive passes
($p=0.5682$). Pro versus Luna has 18 versus five exclusive passes
($p=0.0106$, unadjusted). We therefore describe the observed range and paired
differences rather than asserting that each successive backend is better.

The strongest configuration still fails 35 frozen tasks. Twenty-eight tasks
are passed by none of the six backends, while 17 are passed by all six. Potential
contract defects contribute to these raw counts; Section~\ref{sec:defects}
shows the corresponding exclusion sensitivity. The main result is substantial
but incomplete success under the recorded model--harness conditions, rather
than a defect-free estimate of all software-reuse capability.

\subsection{RQ2: Evaluator Failure Outcomes}
\label{sec:rq2}

The exclusive outcomes in Table~\ref{tab:funnel} separate failed delivery from
failures of the collected artifact. Pro and Flash produce a package for every
task and have no Build-first failures. Of their 77 combined failures, 76 first
fail at Public or Hidden and one at Isolation. Weaker configurations add both
delivery and loading failures. Qwen has 25 empty submissions; GLM has 38. These
remain functional failures but are not semantic diagnoses of code that was
never delivered.

\begin{table}[t]
 \caption{Mutually exclusive first outcomes; each row sums to 150.}
 \label{tab:funnel}
 \centering
 \begin{tabular}{@{}lrrrrrr@{}}
 \toprule
 Backend & Pass & Missing & Build & Public & Hidden & Isolation \\
 \midrule
 Pro & 115 & 0 & 0 & 25 & 10 & 0 \\
 Flash & 108 & 0 & 0 & 26 & 15 & 1 \\
 Luna & 102 & 6 & 3 & 27 & 12 & 0 \\
 GLM & 68 & 38 & 3 & 31 & 8 & 2 \\
 Qwen & 63 & 25 & 6 & 35 & 21 & 0 \\
 OSS & 36 & 2 & 18 & 70 & 23 & 1 \\
 \bottomrule
 \end{tabular}
\end{table}

The non-exclusive view provides a different perspective
(Table~\ref{tab:independent-gates}). Among 829 delivered artifacts, 244 (29.4\%)
have a failed Public flag and 296 (35.7\%) a failed Hidden flag, compared with
39 (4.7\%) failed Isolation flags. Only four artifacts pass the other three
gates and fail solely at Isolation. This supports a narrow operational finding:
the separate Isolation gate rarely determines the final result once the other
gates pass. It does not show that source independence is easy or that packaging
mistakes are absent; source-free execution can expose those mistakes earlier.

\begin{table}[t]
 \caption{Non-exclusive failed gate flags among delivered artifacts. Counts
 can overlap. Iso. residual requires the other three gates to pass.}
 \label{tab:independent-gates}
 \centering\small
 \begin{tabular}{@{}lrrrrrr@{}}
 \toprule
 Backend & Delivered & Build & Public & Hidden & Isolation & Iso. residual \\
 \midrule
 Pro & 150 & 0 & 25 & 29 & 0 & 0 \\
 Flash & 150 & 0 & 26 & 36 & 1 & 1 \\
 Luna & 144 & 3 & 30 & 35 & 3 & 0 \\
 GLM & 112 & 3 & 34 & 38 & 5 & 2 \\
 Qwen & 125 & 6 & 41 & 58 & 8 & 0 \\
 OSS & 148 & 18 & 88 & 100 & 22 & 1 \\
 \midrule
 Total & 829 & 30 & 244 & 296 & 39 & 4 \\
 \bottomrule
 \end{tabular}
\end{table}

\subsection{RQ3: Observed Task Difficulty}
\label{sec:rq3}

Figure~\ref{fig:difficulty} shows the number of evaluated backends passing each
task. Twenty-two of the 28 all-fail tasks belong to hard3. Pro drops from 94/100
on Core to 21/50 on hard3; Flash drops from 90/100 to 18/50, and Luna from
82/100 to 20/50. GLM and Qwen show the same direction. OSS is near 25\% in
both groups (25/100 versus 11/50), so the contrast is not uniform across all
levels of observed capability.

\begin{figure}[t]
 \centering
 \includegraphics[width=0.92\linewidth]{task_difficulty.pdf}
 \caption{Solve-frequency spectrum for the frozen 150 tasks. The horizontal
 axis counts passing backends out of six. Construction labels describe the
 sampled tasks; they do not define the observed outcome.}
 \Description{A stacked bar chart groups tasks by how many of six backends pass.
 The zero-pass group contains 6 Core and 22 hard3 tasks; the six-pass group
 contains 15 Core and 2 hard3 tasks.}
 \label{fig:difficulty}
\end{figure}

Lift type and construction group are strongly associated: 53 of 56 Direct
tasks belong to Core, while 15 of 18 Composite tasks belong to hard3. A raw
Direct--Adapted--Composite success gradient therefore cannot be attributed to
lift type alone. In the six-backend logistic analysis, hard3 is associated
with lower success (odds ratio 0.149, $p=1.07\times10^{-7}$). Composite versus
Direct is not distinguishable at conventional thresholds after adjustment
(odds ratio 0.655, $p=0.486$). Restricting to Pro, Flash, and Luna gives the
same qualitative pattern: hard3 odds ratio 0.092, while Composite has odds
ratio 0.654 and $p=0.554$.

These results support the construction group as a useful empirical stratifier
within Python-150. They do not establish that the group label causes difficulty,
or that lift type has no effect in a larger, differently balanced sample.
The small Composite group and its concentration in hard3 limit separation of
these factors. We do not infer difficulty from the additional Hard-50 extension
or from the uniformly high entanglement-level label.

\subsection{RQ4: Extraction Footprint on Common Successes}
\label{sec:rq4}

Table~\ref{tab:compactness} describes the passing artifacts. Pro and Flash have
median RRES near one and median detected copy fractions above 0.95. Luna has
median RRES 0.815 and median copy fraction 0.164. These medians involve different
passing sets and cannot alone establish a paired difference. They also conceal
substantial within-backend variation: Luna's copy-fraction interquartile range
is [0.048, 0.967], while Pro's is [0.716, 0.987].

\begin{table}[t]
 \caption{Descriptive statistics on each backend's passing artifacts. Different
 rows have different task sets; lower copying is not inherently better.}
 \label{tab:compactness}
 \centering
 \begin{tabular}{@{}lrrrr@{}}
 \toprule
 Backend & Passes & Mean RRES & Median RRES & Median copy fraction \\
 \midrule
 Pro & 115 & 1.781 & 0.993 & 0.957 \\
 Flash & 108 & 1.730 & 0.998 & 0.968 \\
 Luna & 102 & 0.953 & 0.815 & 0.164 \\
 GLM & 68 & 2.473 & 1.013 & 0.947 \\
 Qwen & 63 & 1.556 & 0.975 & 0.757 \\
 OSS & 36 & 1.176 & 0.983 & 0.180 \\
 \bottomrule
 \end{tabular}
\end{table}

On the 97 tasks passed by both Pro and Luna, median RRES is 0.993 versus
0.697, and median copy fraction is 0.966 versus 0.191
(Figure~\ref{fig:paired-copy}). Pro has higher detected copying on 76 tasks,
Luna on 18, with three ties. The paired Wilcoxon test yields
$p=1.18\times10^{-13}$ and matched-pairs rank-biserial effect size $r=0.881$.
The Flash--Luna comparison on 92 common passes has median RRES 0.999 versus
0.713 and copy fraction 0.970 versus 0.204
($p=4.24\times10^{-13}$ for copying, $r=0.885$).

\begin{figure}[t]
 \centering
 \includegraphics[width=\linewidth]{paired_copy.pdf}
 \caption{Pro versus Luna on 97 common passing tasks: paired copy fractions,
 empirical cumulative distributions, and box plots. The statistic measures
 detected source overlap; footprint is measured separately by RRES.}
 \Description{A scatter plot, empirical cumulative distribution, and box plots
 compare detected copy fractions on the same tasks. Pro has median 0.966 and
 Luna 0.191, with substantial within-model variation.}
 \label{fig:paired-copy}
\end{figure}

Functional success thus coexists with systematic differences in artifact size
and source overlap even after fixing the task. A high copy fraction can describe
a faithful small extraction, and a low fraction can describe a large rewrite;
the two measures should remain distinct. These results do not establish that
Luna is globally better, that copying is undesirable, or that smaller artifacts
are easier to maintain. They show why a pass score alone does not characterize
the delivered reusable artifact.

\section{Exploratory Diagnostics and Sensitivity}
\label{sec:diagnostics}

\subsection{Contract Defects and Denominators}
\label{sec:defects}

Failure inspection flagged seven tasks for potential disagreement between the
public contract and evaluator expectations. The issues include omitted API
members and conflicting expectations about whitespace, literal markers, or
default version formatting. The artifact labels these cases
\texttt{benchmark\_invalid\_candidate}; the review is AI-assisted rather than
an independent human adjudication. Their identities and the exclusion
calculation are recorded in Appendix~\ref{app:defects}.

The main tables retain all 150 tasks and their recorded outcomes. The mechanism
analysis instead removes 14 affected Pro/Flash failure rows, leaving 63 rows
from the original 77. These are different denominators for different questions.
To expose the effect of the flagged tasks on measured capability, we also remove
all seven task IDs from every backend's existing results. This conservative
task-level exclusion does not depend on whether a particular backend passed a
flagged task. It produces a 143-task sensitivity view, not a new frozen suite
or an estimate corrected for all possible defects.

% SENSITIVITY_INSERT

\subsection{Behavioral-Contract Recovery}

The existing close-read examines the public contract, first-failure evidence,
and submitted code for every Pro/Flash artifact failure, with trajectory
inspection to check source access. After the defect-row exclusion, the common
symptoms are behavioral drift and incomplete API obligations. We use
\emph{contract-closure gap} as a descriptive term for a delivered artifact
that omits or changes required observable behavior despite implementing a
substantial part of the target. A Hidden failure alone is insufficient for this
label: the failed expectation must first be supported by the public contract.

For example, in an Alembic revision-map task, both backends expose the requested
lookup API but resolve a literal revision identifier as a symbolic base marker.
The failure concerns precedence between two interpretations, rather than an
absent top-level function. An aiohttp task similarly contains validation logic
while leaving a required exceptional path incomplete. These examples illustrate
how apparent implementation coverage can coexist with a missing behavioral
obligation; they do not disclose protected test assertions.

Every retained Pro/Flash failure has recorded repository inspection. This rules
out a narrow explanation that the agent never inspected source files. It does
not establish that all relevant implementation regions were located or understood.
The coding therefore concerns observable artifact symptoms, not a causal
decomposition of search, reasoning, and editing. The analysis remains an
assistant first pass, without independent human agreement measurements.
Appendix~\ref{app:diagnostic-counts} provides provisional counts and a separate
stratified sample from Luna, Qwen, and OSS; neither is presented as a validated
cross-model root-cause distribution.

\subsection{Delivery and Interaction Diagnostics}

Process evidence helps interpret failed delivery. Qwen's 25 empty Python-150
submissions are associated with tool-validation failures involving required
tool-call fields. GLM's 38 empty submissions are predominantly early terminations
without a package; only one is tagged as a tool-validation event. These are
different model--harness outcomes and should not be pooled into behavioral
failure categories. Conversely, 82 Flash passing artifacts have runner status
other than passed, illustrating why workflow status cannot replace evaluation.

Provider token accounting is incomplete and heterogeneous. Luna and GLM lack
verified token totals for this analysis. Pro and Flash expose cache accounting,
allowing an incremental measure of uncached prompt plus completion tokens;
their overall medians are approximately 88.0k and 123.6k. Qwen and OSS have
provider total-token counts without the same cache breakdown. We therefore
do not rank all six backends by token cost.

Within Pro and Flash, hard3 is associated with more interaction and lower
success. Median incremental tokens rise from 81.5k to 118.3k for Pro and from
108.3k to 202.9k for Flash between Core and hard3. This is an observational
association: difficult tasks may induce longer attempts. It does not show
that additional budget is ineffective, or that a different stopping policy
would improve performance.

\section{Discussion}
\label{sec:discussion}

\paragraph{A boundary-sensitive evaluation of software reuse.}
The benchmark makes the output boundary part of correctness. The upstream
project provides implementation evidence, but the submitted artifact must work
without that project. This links the long-standing reuse and transplantation
problem to the capabilities of interactive coding agents. The contribution is
the task collection and measurement protocol, rather than a claim that software
extraction is a newly discovered activity.

\paragraph{From visible APIs to complete behavior.}
Strong backends frequently produce loadable packages while missing required
behavior. The exploratory cases suggest that tracking API presence alone can
overstate implementation completeness: a branch, default, state transition, or
exception can remain incorrect behind an otherwise plausible interface. A useful
direction is to connect each public obligation to repository evidence and
implementation evidence. The present study motivates that direction but does
not demonstrate that an obligation ledger, additional reflection, or a repair
loop improves the benchmark score.

\paragraph{Preserve multiple dimensions of extraction quality.}
A correct artifact can be large relative to its reference, while a small
artifact can fail required behavior. Common-success comparisons isolate
footprint differences from differing task success sets, but still do not
measure downstream maintenance effort. Reporting correctness, relative size,
and detected copying separately avoids rewarding an incomplete implementation
for being small or penalizing faithful reuse solely for retaining source lines.

\paragraph{Benchmark validation is an ongoing evidence problem.}
Stable reference execution and mechanically consistent metadata are necessary
checks, but do not establish that every evaluator expectation follows the task
statement. The later defect flags make that limitation concrete. Freezing the
recorded results, disclosing suspect tasks, and labeling exclusion sensitivity
allows readers to assess their influence without changing the benchmark
retrospectively. Future releases can repair adjudicated defects under a new
identity; the current results remain attached to their original freeze.

\section{Threats to Validity}
\label{sec:threats}

\paragraph{Construct validity.}
Finite behavioral tests cannot prove equivalence to every behavior of an upstream
feature. The output contract deliberately covers a bounded capability, and
passing should be interpreted within that scope. Contract/evaluator disagreements
can misclassify agents; the seven flagged tasks and sensitivity analysis expose
known concerns without certifying the remaining tasks. Build allows path loading
and therefore does not prove wheel validity or clean installability. Isolation
is also an operational check: dependency mistakes may surface at other gates,
and the environment's available packages can mask undeclared dependencies.

\paragraph{Internal validity.}
Outcomes depend on source materialization, dependency locks, agent/evaluator
images, prompt and context handling, model endpoints, and recovery logic.
The release records identities for these components, and the paper does not pool
earlier freezes or information ablations into the main comparison. Nevertheless,
provider behavior and model-specific tool-call compatibility remain part of the
measured system. One final collected outcome per task does not remove effects
of preflight replacement or transient recovery. Reproduction must preserve the
recorded campaign policy as well as the visible task specification.

\paragraph{External and statistical validity.}
The task set emphasizes Python libraries and developer tooling. It does not
establish results for other languages, interactive applications, distributed
services, or industrial migration workloads. Popular source projects may have
appeared in training data; source availability is intentional, while training
contamination is not measured. Construction-group and lift-type associations
are specific to a curated, imbalanced collection. A single outcome per cell
does not estimate sampling variation across repeated runs. Reported task-level
intervals and exploratory tests also do not fully account for shared repository
origins or multiple comparisons.

\paragraph{Annotation and footprint validity.}
Failure and defect labels are AI-assisted and have no independent human L2
validation. Their scope is a strong-backend census and a separate stratified
sample, not a representative six-backend causal taxonomy. Repository access is
not proof of successful localization. RRES depends on the feasible reference
and line-counting convention, and direct-copy detection is a syntactic proxy.
Neither metric establishes minimality, maintainability, or semantic superiority.
The post hoc defect exclusion is derived from observed failures and is not an
exhaustive fairness audit.

\section{Conclusion}

\flb evaluates whether a coding agent can recover a declared capability from a
real repository and deliver an artifact that works without the source project.
The frozen release contains 200 tasks; a common 150-task subset supports the
six-backend study. Under the recorded OpenHands Main protocol, functional pass
rates span 24.0\%--76.7\%. Behavioral evaluation accounts for most strong-backend
failures, and the hard3 construction group is associated with lower success.
Common passing tasks also reveal distinct extraction footprints and source
overlap. These results support evaluating behavior preservation and artifact
construction as separate dimensions of repository-level software reuse.
Exploratory failure evidence and disclosed task-defect sensitivity qualify the
interpretation rather than replacing the frozen measured outcomes.

\FloatBarrier
\bibliographystyle{ACM-Reference-Format}
\bibliography{references}
\FloatBarrier
\appendix

\section{Supplementary Release Results}
\label{app:extension}

Table~\ref{tab:extension} reports the five available 200-task campaigns under
freeze v2, including their Python-150 and additional Hard-50 strata. The
preflight-replacement tasks are already included in these totals. Pro has no
corresponding 200-task campaign in this record. The additional Hard-50 group is
not the hard3 group inside the main 150 tasks. Flash's 49/50 on this extension
also illustrates why the historical group name should not be read as a universal
empirical difficulty label. We use the extension for release transparency,
not for a six-backend difficulty or compactness comparison.

\begin{table}[h]
 \caption{Supplementary functional results on the 200-task freeze v2 release.}
 \label{tab:extension}
 \centering
 \begin{tabular}{@{}lrrr@{}}
 \toprule
 Backend & Python-150 & Additional Hard-50 & Combined \\
 \midrule
 Flash & 108/150 & 49/50 & 157/200 (78.5\%) \\
 Luna & 102/150 & 42/50 & 144/200 (72.0\%) \\
 GLM & 68/150 & 30/50 & 98/200 (49.0\%) \\
 Qwen & 63/150 & 23/50 & 86/200 (43.0\%) \\
 OSS & 36/150 & 25/50 & 61/200 (30.5\%) \\
 \bottomrule
 \end{tabular}
\end{table}

\section{Flagged Task Identities and Sensitivity}
\label{app:defects}

The existing first-pass defect summary identifies the following task IDs:
\begin{itemize}
 \item \texttt{click\_\_lazy\_command\_core\_\_hard3\_001};
 \item \texttt{flake8\_\_plugin\_options\_core\_\_hard3\_001};
 \item \texttt{hatch\_\_project\_metadata\_core\_\_hard3\_001};
 \item \texttt{pluggy\_\_hook\_wrapper\_core\_\_hard3\_001};
 \item \texttt{pytest\_\_ini\_markers\_core\_\_001};
 \item \texttt{readme\_renderer\_\_content\_type\_core\_\_hard3\_001}; and
 \item \texttt{setuptools\_scm\_\_version\_normalize\_core\_\_hard3\_001}.
\end{itemize}
The sensitivity script removes these same IDs for every backend and recomputes
counts from the existing 900-row result matrix. It validates key uniqueness,
common task membership, and agreement with the recorded evaluator gates.
Exclusion does not assert that every failure on these tasks has the same cause:
one backend may fail earlier for an unrelated dependency problem. These task
flags are provisional and were discovered through failure analysis, which
limits any claim that the remaining subset is free of defects.

\section{Provisional Failure Coding}
\label{app:diagnostic-counts}

Table~\ref{tab:semantic} records the existing AI-assisted L1 coding after
removing 14 suspected defect rows from 77 Pro/Flash artifact failures. No empty
submissions enter this denominator. The analysis contains 63 rows: 28 Pro and
35 Flash. The 54 drift and seven API-completion rows motivate the qualitative
contract-closure discussion. These counts are an auditable first-pass coding
result, not independently validated root-cause proportions.

\begin{table}[h]
 \caption{Provisional L1 primary symptoms on Pro/Flash failures. This is an
 AI-assisted diagnostic census, not independent human gold.}
 \label{tab:semantic}
 \centering
 \begin{tabular}{@{}lrrr@{}}
 \toprule
 Symptom & Pro & Flash & Total \\
 \midrule
 Behavioral drift & 26 & 28 & 54 \\
 Contract/API completion & 2 & 5 & 7 \\
 Packaging/modularization & 0 & 2 & 2 \\
 \midrule
 Total & 28 & 35 & 63 \\
 \bottomrule
 \end{tabular}
\end{table}

A separate stratified sample reviews 45 artifact failures from Luna, Qwen,
and OSS. Thirteen rows are excluded as suspected task defects, leaving 32:
14 drift, nine API-completion, four dependency-closure, and five
packaging/modularization cases. These counts are not pooled with the Pro/Flash
census because the sampling and backend scope differ. GLM's artifact failures
were not included in this semantic sample. Zero primary localization labels in
these inspections do not establish that localization is solved.

\section{Artifact Identity and Reproduction}
\label{app:reproduction}

The active release is freeze v2 with candidate prefix \texttt{212930ea}.
Its full freeze identity is
\begin{quote}\small\ttfamily
6c20ff0307762503a73cbb9ff32e9992c6446\\
e4b17483a68373027be58cbf419
\end{quote}
% The split representation above is typographic only; concatenate without spaces.
The machine-readable manifest retains the unbroken identifier, task-specific
specification and source hashes, and agent/evaluator image digests. The evaluator
image is \texttt{featureliftbench-eval:python200-prime-212930ea}; the agent image
uses the corresponding agent tag. The freeze records Python 3.11.14 and
OpenHands CLI 1.16.0.

The main quantitative input is the existing
\texttt{python150\_prime\_v2\_analysis\_20260905/task\_results.csv} analysis
matrix. The final analysis package retains summary tables, paired-task records,
statistical outputs, and figures. Raw campaign directories retain per-task
submissions, evaluator results, run metadata, trajectories, and usage records.
The paper revision adds only a reproducible offline known-defect sensitivity
calculation; it does not alter task definitions or original result files.

Upstream snapshots retain their original licenses. The distribution separates
source acquisition, evaluator-side assets, and submitted artifacts, whose
licensing obligations must be respected individually. The local evidence bundle
supports the results in this draft; a public archival identifier and finalized
distribution metadata are not established here. Earlier freezes, development
pilots, alternative information arms, and unrun method proposals are outside the
main evidence set.

\end{document}
